\section{The large-concentration limit of the direct jump}\label{app:hdp}

In this appendix we prove Proposition~\ref{prop:directjump}(d). As in the
proof of part~(c), $B_1,B_2,B_3$ are independent with the law
$\mathrm{Beta}(a,a)$, where $a=\omega/3$ throughout, and $\Prob_\omega(D)$ is the
probability of the event $D$ defined there. We let $a\to\infty$, so we may
assume $a\ge1$.

For large $a$ each $B_j$ lies within about $a^{-1/2}$ of $\frac12$. At the
uniform kernel $B_1=B_2=B_3=\frac12$ all 3 constraints are equalities. To
first order they allow only the cyclic kernels of
Proposition~\ref{prop:directjump}(b), so $D$ is a thin tube around this
curve of kernels. Its cross-section at distance $\delta$ from the uniform
kernel has area of order $\delta^4$. The density of $(B_1,B_2,B_3)$ near the
uniform kernel is of order $a^{3/2}$, so $\Prob_\omega(D)$ is of order
$a^{3/2}\int\delta^4e^{-ca\delta^2}\,d\delta$, that is of order $1/a$.

\paragraph{Logits.}
Put $\ell_j=\log\bigl(B_j/(1-B_j)\bigr)$ and $g(y)=2\log\cosh(y/2)$. The
function $g$ is even, convex and nonnegative, with $g(0)=g'(0)=0$. Since
$B_j=1/(1+e^{-\ell_j})$, we have
$2B_j=e^{\ell_j/2}/\cosh(\ell_j/2)$,
$2(1-B_j)=e^{-\ell_j/2}/\cosh(\ell_j/2)$ and
$B_j(1-B_j)=\frac14e^{-g(\ell_j)}$. As $dB_j/d\ell_j=B_j(1-B_j)$, the logit
$\ell_j$ has the density
$\bigl(B_j(1-B_j)\bigr)^a\,\Gamma(2a)/\Gamma(a)^2=e^{-ag(\ell_j)}/Z$ with
$Z=4^a\,\Gamma(a)^2/\Gamma(2a)=2\sqrt\pi\,\Gamma(a)/\Gamma(a+\frac12)$ by the
duplication formula. Write $g_j=g(\ell_j)$. Then
$\log(2B_j)=\frac12(\ell_j-g_j)$ and $\log(2(1-B_j))=\frac12(-\ell_j-g_j)$.
So after taking logarithms the 3 constraints of $D$ read
\begin{equation}\label{eq:hdp-constraints}
 M_1\le g_1+g_2,\qquad M_2\le g_1+g_3,\qquad M_3\le g_2+g_3,
\end{equation}
where $M_1=\ell_1+\ell_2$, $M_2=\ell_3-\ell_1$ and $M_3=-\ell_2-\ell_3$. Note
that $M_1+M_2+M_3=0$. From now on we regard $D$ as the set of
$\ell\in\mathbb R^3$ that satisfy \eqref{eq:hdp-constraints}. With
$\psi(\ell)=g_1+g_2+g_3$ we get
\begin{equation}\label{eq:hdp-integral}
 \Prob_\omega(D)=Z^{-3}\int_D e^{-a\psi(\ell)}\,d\ell .
\end{equation}

\paragraph{Two bounds on $g$.}
For all real $y$,
\begin{align}
 \frac{y^2}4-\frac{y^4}{96}\le g(y)&\le\frac{y^2}4,
 \tag{g1}\label{eq:hdp-g1}\\
 g(y)&\ge\frac1{10}\min\bigl(y^2,\lvert y\rvert\bigr).
 \tag{g2}\label{eq:hdp-g2}
\end{align}
Put $u=y/2$, so that $g(y)=2\log\cosh u$. The upper bound in
\eqref{eq:hdp-g1} is $\cosh u\le e^{u^2/2}$, which holds term by term in the
power series since $(2m)!\ge2^mm!$. For the lower bound let
$\chi(u)=\log\cosh u-\frac{u^2}2+\frac{u^4}{12}$. Then $\chi(0)=\chi'(0)=0$
and $\chi''(u)=u^2-\tanh^2u\ge0$, so $\chi\ge0$. Since
$2(\frac{u^2}2-\frac{u^4}{12})=\frac{y^2}4-\frac{y^4}{96}$, this is the lower
bound in \eqref{eq:hdp-g1}. For \eqref{eq:hdp-g2}, note that
$g''(y)=\frac12\operatorname{sech}^2(y/2)$ is at least
$\frac12\operatorname{sech}^21>0.2$ on $\lvert y\rvert\le2$. With
$g(0)=g'(0)=0$ this gives $g(y)\ge0.1y^2$ there. For $\lvert y\rvert\ge2$,
convexity and $g(0)=0$ give $g(y)\ge\frac12\lvert y\rvert g(2)$, and
$\frac12g(2)=\log\cosh1>0.43$.

\paragraph{Coordinates.}
Let $v=(1,-1,1)$ and $t=\langle\ell,v\rangle/3$. The linear map
$\ell\mapsto(t,M_1,M_2)$ has determinant $\frac13+\frac13+\frac13=1$, as we
see by expanding along the row of $t$. Its inverse is $\ell=tv+w$ with
$w=\frac13(M_1-M_2,\,2M_1+M_2,\,M_1+2M_2)$, and $w$ is orthogonal to $v$.
Hence $\lvert\ell\rvert^2=3t^2+\lvert w\rvert^2$, and a direct computation
gives $\lvert w\rvert^2=\frac13(M_1^2+M_2^2+M_3^2)$. We write $\lvert M\rvert$ for
the Euclidean norm of $(M_1,M_2)$ and
$\lvert M\rvert_\infty=\max(\lvert M_1\rvert,\lvert M_2\rvert)$. Since
$M_3=-(M_1+M_2)$,
\begin{equation}\label{eq:hdp-w}
 \tfrac13\lvert M\rvert^2\le\lvert w\rvert^2\le\lvert M\rvert^2
 \le2\lvert M\rvert_\infty^2 .
\end{equation}
The line $M=0$ is $\ell=tv$, that is $B_1=B_3=1-B_2$. These are the cyclic
kernels $\pi_{i,i+1}=B_1$ of Proposition~\ref{prop:directjump}(b).

\begin{proof}[Proof of Proposition~\ref{prop:directjump}\textup{(d)}]
\emph{Step 1.} We bound $M$ on $D$. By \eqref{eq:hdp-constraints} and
$g\ge0$ we have $M_1\le g_1+g_2\le\psi$ and
$M_1=-M_2-M_3\ge-(g_1+g_3)-(g_2+g_3)\ge-2\psi$. The same argument applies to
$M_2$. With \eqref{eq:hdp-g1} we get, on $D$,
\begin{equation}\label{eq:hdp-apriori}
 \lvert M\rvert_\infty\le2\psi(\ell)\le\frac{\lvert\ell\rvert^2}2
 =\frac{3t^2+\lvert w\rvert^2}2 .
\end{equation}

\emph{Step 2.} We show that the far region is negligible. Fix
$\varepsilon=\frac1{12}$ and let
$U=\{\lvert t\rvert\le\varepsilon,\ \lvert M\rvert_\infty\le\varepsilon\}$.
Off $U$ we have $\lvert\ell\rvert\ge\varepsilon/\sqrt3$. Indeed, either
$\lvert t\rvert>\varepsilon$ and $\lvert\ell\rvert\ge\sqrt3\lvert t\rvert$, or
$\lvert M\rvert_\infty>\varepsilon$ and
$\lvert\ell\rvert\ge\lvert w\rvert\ge\lvert M\rvert/\sqrt3$ by
\eqref{eq:hdp-w}. Some $\lvert\ell_j\rvert$ is at least
$\lvert\ell\rvert/\sqrt3$, and $\min(y^2,\lvert y\rvert)$ increases with
$\lvert y\rvert$. So \eqref{eq:hdp-g2} gives $\psi(\ell)\ge\gamma(\ell)$ with
$\gamma(\ell)=0.1\min\bigl(\lvert\ell\rvert^2/3,\lvert\ell\rvert/\sqrt3\bigr)$.
As $\gamma$ increases with $\lvert\ell\rvert$, it is at least
$\gamma_0=\varepsilon^2/90$ off $U$. Hence
$e^{-a\psi}\le e^{-(a-1)\gamma_0}e^{-\gamma}$ off $U$. Since $\gamma$ grows
linearly, $e^{-\gamma}$ is integrable over $\mathbb R^3$, and so
$\int_{D\setminus U}e^{-a\psi(\ell)}\,d\ell=O\bigl(e^{-\gamma_0a}\bigr)$.

\emph{Step 3.} We show that the sections of $D\cap U$ are close to
triangles. Let $(t,M)\in D\cap U$. By \eqref{eq:hdp-apriori} and
\eqref{eq:hdp-w},
$\lvert M\rvert_\infty\le\frac32t^2+\lvert M\rvert_\infty^2
\le\frac32t^2+\varepsilon\lvert M\rvert_\infty$, so
$\lvert M\rvert_\infty\le2t^2$.

Next let $E$ be the set of $(t,M)$ with $\lvert t\rvert\le\varepsilon$ and
$\lvert M\rvert_\infty\le2t^2$. On $E$, \eqref{eq:hdp-w} gives
$\lvert w\rvert\le3t^2$, so $\ell_j=\pm t+w_j$ with
$\lvert w_j\rvert\le3t^2$. As $\lvert t\rvert\le\frac1{12}$, this gives
$\lvert\ell_j\rvert\le2\lvert t\rvert$ and
$\lvert\ell_j^2-t^2\rvert\le6\lvert t\rvert^3+9t^4\le7\lvert t\rvert^3$. Then
\eqref{eq:hdp-g1} gives
$\lvert g_j-\frac14t^2\rvert\le\frac1{96}\ell_j^4+\frac74\lvert t\rvert^3
\le2\lvert t\rvert^3$. So on $E$
\begin{equation}\label{eq:hdp-pair}
 \bigl\lvert g_i+g_j-\tfrac12t^2\bigr\rvert\le4\lvert t\rvert^3
 \qquad(i\ne j).
\end{equation}
Also on $E$, \eqref{eq:hdp-g1} gives
$\frac14\lvert\ell\rvert^2-\frac12t^4\le\psi(\ell)\le\frac14\lvert\ell\rvert^2$,
and $\lvert\ell\rvert^2=3t^2+\lvert w\rvert^2$ with
$\lvert w\rvert^2\le8t^4$. Hence
\begin{equation}\label{eq:hdp-psi}
 \bigl\lvert\psi(\ell)-\tfrac34t^2\bigr\rvert\le2t^4 .
\end{equation}

For $h\ge0$ let
$\mathcal T(h)=\{M:M_1\le h,\ M_2\le h,\ M_1+M_2\ge-h\}$. This is the
triangle with vertices $(h,h)$, $(h,-2h)$ and $(-2h,h)$. It has area
$\frac92h^2$ and lies in $\{\lvert M\rvert_\infty\le2h\}$. For
$\lvert t\rvert\le\varepsilon$ let $D_t$ be the set of $M$ with
$(t,M)\in D\cap U$, and put $h_\pm=\frac12t^2\pm4\lvert t\rvert^3$. Note that
$h_-\ge0$, since $8\lvert t\rvert<1$. We claim that
$\mathcal T(h_-)\subseteq D_t\subseteq\mathcal T(h_+)$. If $M\in D_t$, then
$(t,M)\in E$ by the first part of this step. So \eqref{eq:hdp-constraints}
and \eqref{eq:hdp-pair} give $M_1\le h_+$, $M_2\le h_+$ and
$-(M_1+M_2)=M_3\le h_+$. If $M\in\mathcal T(h_-)$, then
$\lvert M\rvert_\infty\le2h_-\le t^2$, and $t^2\le\min(2t^2,\varepsilon)$.
So $(t,M)\in E\cap U$. Each of $M_1,M_2,M_3$ is at most $h_-$, so
\eqref{eq:hdp-pair} gives \eqref{eq:hdp-constraints}, and $M\in D_t$. Since
$h_\pm=\frac12t^2(1\pm8\lvert t\rvert)$, the area of $D_t$ satisfies
\begin{equation}\label{eq:hdp-area}
 \tfrac98t^4(1-8\lvert t\rvert)^2\le\lvert D_t\rvert
 \le\tfrac98t^4(1+8\lvert t\rvert)^2 .
\end{equation}

\emph{Step 4.} We apply Laplace's method. The coordinates $(t,M_1,M_2)$ have
Jacobian $1$. So by Fubini's theorem, \eqref{eq:hdp-psi} and
\eqref{eq:hdp-area}, the integral $J=\int_{D\cap U}e^{-a\psi}\,d\ell$
satisfies
\[
 \int_{-\varepsilon}^{\varepsilon}\tfrac98t^4(1-8\lvert t\rvert)^2
 e^{-a(\frac34t^2+2t^4)}\,dt
 \le J\le
 \int_{-\varepsilon}^{\varepsilon}\tfrac98t^4(1+8\lvert t\rvert)^2
 e^{-a(\frac34t^2-2t^4)}\,dt .
\]
Substitute $t=s/\sqrt a$. Both bounds become $a^{-5/2}$ times an integral
over $\lvert s\rvert\le\varepsilon\sqrt a$, and both integrands tend to
$\frac98s^4e^{-3s^2/4}$ for each $s$. On $\lvert s\rvert\le\varepsilon\sqrt a$
we have $2s^4/a\le2\varepsilon^2s^2\le s^2/4$ and
$8\lvert s\rvert/\sqrt a\le8\varepsilon<1$. So both integrands are at most
$\frac92s^4e^{-s^2/2}$, which is integrable. By dominated convergence and
$\int_{\mathbb R}s^4e^{-bs^2}\,ds=\frac34\sqrt\pi\,b^{-5/2}$ with
$b=\frac34$,
\[
 J=a^{-5/2}\Bigl(\frac98\cdot\frac{8\sqrt\pi}{3\sqrt3}+o(1)\Bigr)
 =a^{-5/2}\bigl(\sqrt{3\pi}+o(1)\bigr).
\]
By Step 2 the same holds for the integral over all of $D$.

It remains to estimate $Z$. The Gamma function is log-convex. So
$\Gamma(a+\frac12)^2\le\Gamma(a)\Gamma(a+1)=a\Gamma(a)^2$ and
$\Gamma(a+1)^2\le\Gamma(a+\frac12)\Gamma(a+\frac32)
=(a+\frac12)\Gamma(a+\frac12)^2$. Hence
$a(a+\frac12)^{-1/2}\le\Gamma(a+\frac12)/\Gamma(a)\le\sqrt a$, and
$Z=2\sqrt{\pi/a}\,\bigl(1+O(1/a)\bigr)$. Now \eqref{eq:hdp-integral} gives
\[
 \Prob_\omega(D)=\frac{a^{3/2}}{8\pi^{3/2}}\,a^{-5/2}\sqrt{3\pi}\,\bigl(1+o(1)\bigr)
 =\frac{\sqrt3}{8\pi a}\bigl(1+o(1)\bigr)
 =\frac{3\sqrt3}{8\pi\omega}\bigl(1+o(1)\bigr). \qedhere
\]
\end{proof}
